\section{Related works}
\label{sec:related_works}

\begin{table*}[t]
  \caption{Description of the attributes within the \texttt{requirements} dictionary.}
  \label{tab:requirements_schema}
  \centering
  \begin{tabular}{p{4cm} p{2cm} p{10cm}}
      \toprule
      Attribute & Data Type & Description \\
      \midrule
      \texttt{id} & String & Unique identifier for the building, corresponding to an entry in the \texttt{massing} list. \\
      \texttt{building\_function} & String & The primary function of the building (e.g., ``Commercial Accommodation''). \\
      \texttt{n\_dwellings} & Float & The total number of living spaces (apartments, rooms, etc.). \\
      \texttt{building\_name} & String & The name of the building, if one exists. \\
      \texttt{n\_floors} & Float & The number of floors above ground. \\
      \texttt{n\_commercial\_spaces} & Integer & The total count of commercial spaces (e.g., offices, retail) within the building. \\
      \texttt{offices} & List[Dict] & A list of specific commercial spaces. Each item contains a \texttt{name} and a \texttt{function} describing the business. \\
      \texttt{public\_spaces} & List[Dict] & A list of spaces with public capacity (e.g., cafes, restaurants). Each item contains \texttt{name}, \texttt{function}, \texttt{indoor\_capacity}, and \texttt{outdoor\_capacity}. \\
      \texttt{floor\_height} & Float & The mean height of a single floor, in meters. \\
      \texttt{usable\_area} & Float & The total usable area of the building, calculated as the sum of the areas of all floors, in square meters. \\
      \bottomrule
  \end{tabular}
\end{table*}

\textbf{Algorithmic and Evolutionary Methods.} A significant body of work focuses on generating massing through parametric modeling and evolutionary algorithms. \cite{1, 2, Gaier_2024} exemplify this approach, using predefined geometric rules and performance criteria (e.g., solar exposure, energy efficiency) to generate and optimize massing forms. While these methods are powerful for navigating a defined solution space towards specific objectives, they lack the flexibility to learn implicit, real-world design patterns from data. Their output is constrained by the initial rule set and requires significant expert knowledge to set up, making them less adaptable to novel contextual or programmatic cues compared to data-driven methods.

\textbf{Agentic and Collaborative AI Approaches.} With the rise of large language models, recent work explores using AI as a collaborative agent in the design process. \cite{3, 4} demonstrate frameworks where multimodal AIs converse with human designers to support early-stage exploration. These systems excel at ideation and leveraging common-sense knowledge but operate in a text-centric manner without being grounded in a large-scale, structured dataset of architectural forms and their real-world constraints. Consequently, they lack the capability for direct, end-to-end massing generation from specified site conditions.

\textbf{Data-Driven Deep Learning Methods.} Several studies have applied deep learning directly to geometric generation. \cite{5} uses a Deep Neural Network to classify and investigate high-rise building types, but does not generate geometry. \cite{6} proposes I2bnet, which generates a BIM massing model from a single image, demonstrating a data-driven approach. Similarly, \cite{7} integrates footprint prediction with massing in a Pittsburgh case study. While these works share our data-driven philosophy, they typically do not condition generation on rich, multi-modal inputs (textual requirements and visual context) simultaneously, don't use Vision-Language Models \cite{vlm}, and they operate on a scale and modality different from our VLM-based formulation.

\textbf{Contextual and Urban-Scale Generation.} Research in this area focuses on generating building geometries that respond to their environment. \cite{7} integrates building footprint prediction with massing in a Pittsburgh case study, creating context-aware massing models, but relies on traditional GIS data and parametric modeling rather than learning from a large-scale dataset. \cite{4} presents a framework combining computational optimization with generative AI for performative design, yet it operates at a conceptual level using text-to-image models and does not produce structured, parseable 3D massing geometry. In contrast, our work introduces a dedicated dataset and an end-to-end VLM framework that directly generates detailed, program-specific 3D massing from structured requirements and visual context, moving beyond conceptual inspiration or parametric constraints to data-driven synthesis.
